% Chapter 18 -- KiCad: from schematic to a two-layer sensor board
\chapter{KiCad: from schematic to a two-layer sensor board}\label{ch:kicad}

The breadboard proved the circuit; the printed circuit board is how
you keep it. This chapter walks the complete path a hobbyist's design
actually travels: schematic symbols wired into a netlist, footprints
assigned, copper placed and poured on a real two-layer stack, design
rules checked, and a Gerber bundle zipped off to a fab that will mail
back green fibreglass. The board we design is small and useful --- a
5\,V-to-3.3\,V rail, a TMP36 temperature sensor through an RC filter,
and a header that mates to the Tang Nano's GPIO --- and every file it
produces is real: the netlist, the DRC-style check, and the Gerbers
themselves are generated by the chapter's own scripts.

% from-zero:begin
\begin{rulebox}[Before you start]
Chapters~1--7 and~17. \emph{Recap:} a PCB turns the schematic into copper traces on a board that a factory can make.
\end{rulebox}
\begin{wins}[After this chapter you can]
\begin{itemize}
\item follow the schematic-to-board pipeline in KiCad;
\item choose correct footprints;
\item lay out a two-layer board;
\item export and check the fabrication files.
\end{itemize}
\end{wins}
\begin{principle}[How to think about this]
A circuit board is a schematic made permanent, so mistakes are costly. Work in a pipeline, check each stage before moving on, and think like the factory that will receive your files. Footprints and design rules are the contract between you and it.
\end{principle}
% from-zero:end

\section{The pipeline}\label{sec:workflow18}

\index{KiCad}\index{PCB}
KiCad's flow is a five-stage pipeline with two spell-checkers built
in. The schematic is where the \emph{meaning} lives (which pins
connect to which nets); footprints translate meaning into copper
geometry; the PCB editor places parts and routes tracks; ERC and DRC
are the two audits --- schematic rules before layout, board rules
before fab; and the Gerbers are the contract a fab can manufacture
without ever seeing your schematic.

\begin{figure}[htbp]
\centering
\begin{tikzpicture}[font=\scriptsize, >=latex]
\foreach \i/\t/\d in {0/{schematic}/{symbols + wires},1/{footprints}/{pin $\to$ package},2/{PCB}/{place, pour, route},3/{DRC}/{rules check},4/{Gerbers}/{fab files}}{
  \pgfmathsetmacro{\x}{0.4+\i*2.3}
  \draw[thick, fill=accentsoft] (\x,3.2) rectangle (\x+1.9,4.4);
  \node[align=center] at (\x+0.95,4.0) {\t};
  \node[align=center, font=\tiny, text=slate] at (\x+0.95,3.55) {\d};
}
\foreach \i in {0,...,3}{
  \pgfmathsetmacro{\x}{2.3+\i*2.3}
  \draw[->, thick] (\x,3.8) -- (\x+0.3,3.8);}
\draw[thick, fill=goodbg] (4.0,1.6) rectangle (7.2,2.6);
\node[align=center] at (5.6,2.1) {ERC (schematic rules)\\--- runs \emph{before} footprints};
\draw[->, dashed, slate] (2.3,3.2) -- (4.0,2.4);
\draw[->, dashed, slate] (7.2,2.1) -- (9.0,3.2);
\node[align=left, text=slate, text width=10cm] at (5.3,0.4)
{The pipeline is one-directional but the loop is not: every stage can
send you back one. ERC at schematic time, DRC before fab --- the two
spell-checks that keep a bad board from becoming a real package in
the mail.};
\end{tikzpicture}

\caption{The KiCad pipeline with its two checkpoints. Regenerate:
\ttfamily figures/ch18/fig\_workflow.tex.}
\label{fig:workflow18}
\end{figure}

\section{The board spec}\label{sec:spec18}

\index{sensor board}
Design starts with a sentence: \emph{a small board that gives the
Tang Nano a temperature it can read}. Requirements:

\begin{itemize}
\item \textbf{Power:} 5\,V from a header pin through an AMS1117-3.3
LDO to a 3.3\,V rail (decoupled both sides, per its datasheet).
\item \textbf{Signal:} TMP36 analog sensor $\to$ 1\,k$\Omega$/100\,nF
lowpass ($f_c\approx1.6$\,kHz) $\to$ header pin.
\item \textbf{Mechanics:} 40$\times$30\,mm, two mounting holes,
all hand-solderable parts (0805 / through-hole).
\item \textbf{Interface:} FILT goes to an FPGA pin configured as a
comparator input --- the honest scope: the FPGA thresholds the analog
level; a real ADC is a different (bigger) board.
\end{itemize}

\begin{figure}[htbp]
\centering
\begin{tikzpicture}[font=\scriptsize, >=latex]
% sensor board schematic (deliverable figure)
% power in + LDO
\node[left] at (0,4.6) {+5\,V (J1)};
\draw[thick] (0.3,4.6) -- (1.6,4.6);
\draw[fill=white] (1.6,4.1) rectangle (3.6,5.1);
\node[align=center] at (2.6,4.6) {U1\\AMS1117-3.3};
\draw[thick] (3.6,4.6) -- (5.2,4.6) node[right]{+3.3\,V};
\fill[accent] (4.4,4.6) circle (0.07);
\draw[thick] (4.4,4.6) -- (4.4,5.4) -- (2.0,5.4);
\draw[fill=white] (1.6,5.4) rectangle (2.4,5.8); \node at (2.0,5.6){C4};
\draw[thick] (2.0,5.4) -- (2.0,5.0);
\draw[thick] (4.4,5.4) -- (6.0,5.4);
\draw[fill=white] (5.6,5.4) rectangle (6.4,5.8); \node at (6.0,5.6){C2};
\draw[thick] (6.0,5.4) -- (6.0,4.6);
% sensor
\draw[fill=white] (7.2,3.6) rectangle (8.8,4.6);
\node[align=center] at (8.0,4.1) {U2\\TMP36};
\draw[thick] (7.2,4.3) -- (6.4,4.3) -- (6.4,4.6) -- (6.4,4.6);
\draw[thick] (7.2,3.9) -- (6.8,3.9) -- (6.8,3.4) -- (6.8,3.0);
\node[below] at (6.8,3.0) {GND};
% RC filter
\draw[thick] (8.8,4.1) -- (9.4,4.1);
\draw[fill=white] (9.4,3.9) rectangle (10.6,4.3); \node at (10.0,4.1){R1 1k};
\draw[thick] (10.6,4.1) -- (11.4,4.1) node[right]{FILT $\to$ J1.4};
\fill[accent] (11.0,4.1) circle (0.07);
\draw[thick] (11.0,4.1) -- (11.0,3.2);
\draw[fill=white] (10.7,2.8) rectangle (11.3,3.2); \node[right, font=\tiny] at (11.3,3.0){C1 100n};
\draw[thick] (11.0,2.8) -- (11.0,2.4) node[below]{GND};
\draw[thick] (9.4,4.1) -- (9.4,3.2) -- (10.0,3.2);
\draw[fill=white] (9.7,2.8) rectangle (10.3,3.2); \node[left, font=\tiny] at (9.7,3.0){R2 10k};
\draw[thick] (10.0,2.8) -- (10.0,2.4) node[below]{GND};
% labels
\node[align=left, text=slate, text width=11.5cm] at (5.8,1.6)
{The board's whole schematic: 5\,V in, regulated to 3.3\,V with
decoupling both sides, the TMP36's 750\,mV at 50\textdegree C
through a 1\,k$\Omega$/100\,nF lowpass ($f_c$ = 1.6\,kHz --- the real
sims) to a header pin the FPGA can threshold. R2 bleeds the filter
node so it isn't floating when the sensor is absent.};
\end{tikzpicture}

\caption{The deliverable: the sensor board's complete schematic.
Regenerate: \ttfamily figures/ch18/fig\_schematic.tex.}
\label{fig:schematic18}
\end{figure}

\begin{example}[The signal chain, verified]
Before footprints exist, the filter gets simulated at the real BOM
values --- \texttt{ch18\_sensor\_filter.cir} runs a real \texttt{.ac}
sweep on the 750\,mV-at-50\textdegree C output through 1\,k$\Omega$
and 100\,nF:

\begin{lstlisting}[caption={Real ngspice output --- the filter's
measured corner matches the 1/(2$\pi$RC) prediction.}]
fc                  =  1.59226e+03
\end{lstlisting}

1.59\,kHz --- within 0.1\% of the 1.591\,kHz theory. A temperature
channel has no business carrying faster content anyway; the corner
is the design.
\end{example}

\begin{example}[The netlist: what the schematic exports]
The schematic's product is the netlist --- components and nets, no
geometry. The board's real excerpt:

\begin{lstlisting}[caption={\ttfamily ch18\_sensor.net excerpt ---
the five nets, verbatim.}]
(net (code "1") (name "+5V")
     (node (ref "J1") (pin "1")) (node (ref "U1") (pin "3"))
     (node (ref "C3") (pin "1")) (node (ref "C4") (pin "1")))
(net (code "2") (name "GND")  ... all returns ...)
(net (code "3") (name "+3V3")
     (node (ref "U1") (pin "1")) (node (ref "U2") (pin "1"))
     (node (ref "J1") (pin "3")) ...)
(net (code "4") (name "SENSE")
     (node (ref "U2") (pin "2")) (node (ref "R1") (pin "1")))
(net (code "5") (name "FILT")
     (node (ref "R1") (pin "2")) (node (ref "R2") (pin "1"))
     (node (ref "J1") (pin "4")))
\end{lstlisting}

Six components' worth of connectivity is the whole design --- five
nets: supply, ground, regulated rail, sensor node, filtered output.
If the netlist is right, every later stage is mechanics.
\end{example}

\begin{checkpoint}
Before footprints, run ERC: it finds the pin a wire missed, the
floating output, and the duplicated reference while they're still
one-line fixes.
\end{checkpoint}

\section{Footprints: value versus body}\label{sec:footprint18}

\index{footprint}
A resistor's \emph{value} lives in the schematic; its \emph{size}
lives in the footprint. Choosing 0805 over 0402 or axial is a
physical decision: hand-solderable by a first-timer, big enough to
rework, small enough to not dominate a 40\,mm board.

\begin{figure}[htbp]
\centering
\begin{tikzpicture}[font=\scriptsize, >=latex]
% footprint panel: same part, three packages
\node[font=\small\bfseries] at (4.6,4.8) {one resistor, three bodies};
\foreach \x/\nm/\s in {0.6/{0402}/{1.0x0.5mm},4.4/{0805}/{2.0x1.25mm},8.2/{axial}/{through-hole}}{
  \draw[thick, fill=graybg] (\x,2.6) rectangle (\x+2.8,4.2);
}
\draw[fill=accentsoft] (1.4,3.3) rectangle (2.0,3.5);
\draw[fill=accentsoft] (5.2,3.2) rectangle (6.4,3.6);
\draw[thick] (8.5,3.4) -- (9.9,3.4);
\draw[fill=accentsoft] (9.0,3.2) rectangle (10.0,3.6);
\draw[thick] (10.0,3.4) -- (11.0,3.4);
\node[align=center] at (2.0,2.9) {0402\\(tweezers)};
\node[align=center] at (5.8,2.9) {0805\\(our pick --- hand-solderable)};
\node[align=center] at (9.7,2.9) {axial\\(through-hole, biggest)};
\node[align=left, text=slate, text width=10.5cm] at (5.3,1.7)
{Footprint $\neq$ value: the schematic says ``1\,k$\Omega$,'' the
footprint says how much board it sits on. We stay at 0805 and
through-hole for everything on this board --- hand-solderable,
hand-fixable, and big enough to read.};
\end{tikzpicture}

\caption{Three packages, same resistor: 0402 (tiny), 0805 (our pick),
axial through-hole (bulky). Regenerate: \ttfamily
figures/ch18/fig\_footprint.tex.}
\label{fig:footprint18}
\end{figure}

The board's footprint assignments: SOT-223 for the regulator,
TO-92 inline for the sensor, 0805 for all passives, and a
2.54\,mm pin header for the connector --- nothing on this board
requires solder paste or a microscope.

\section{The two-layer stack}\label{sec:stack18}

\index{two-layer PCB}
The cheapest board has two copper faces and a 1.6\,mm core. Our
honest split: \textbf{top} carries signals and pads (parts live
here); \textbf{bottom} is one unbroken ground pour --- every ground
pin gets a short drop to the same copper sheet instead of a long
snaking track.

\begin{figure}[htbp]
\centering
\begin{tikzpicture}[font=\scriptsize, >=latex]
% two-layer stack, exploded
\draw[thick, fill=accentsoft] (0,3.8) -- (8.0,3.8) -- (9.0,4.6) -- (1.0,4.6) -- cycle;
\node[above] at (4.5,4.6) {F.Cu --- signals + pads (top)};
\draw[thick, fill=graybg] (0.4,2.8) -- (8.4,2.8) -- (9.4,3.6) -- (1.4,3.6) -- cycle;
\node[below] at (4.9,2.8) {FR4 core, 1.6\,mm};
\draw[thick, fill=goodbg] (0.8,1.8) -- (8.8,1.8) -- (9.8,2.6) -- (1.8,2.6) -- cycle;
\node[below] at (5.3,1.8) {B.Cu --- ground pour (bottom)};
% via
\draw[thick, slate, dashed] (5.8,4.4) -- (5.8,2.1);
\node[right, font=\tiny] at (5.9,3.1) {via};
% mask legend
\draw[thick, fill=white] (9.4,0.6) rectangle (11.6,1.4);
\node[align=left, font=\tiny] at (10.5,1.0) {soldermask over both\\faces (green in photos)};
\node[align=left, text=slate, text width=10.5cm] at (5.3,0.3)
{Two copper layers and a core: signal and parts on top, one big
ground pour underneath --- the cheapest board fabs sell and the
easiest to lay out honestly.};
\end{tikzpicture}

\caption{The stack: signals on F.Cu, ground on B.Cu, FR4 between,
vias crossing. Regenerate: \ttfamily figures/ch18/fig\_stack.tex.}
\label{fig:stack18}
\end{figure}

Placement is where the schematic's intent meets the board's physics:
connector on the edge the cable exits, the LDO near the input with
its capacitors touching its pins (decoupling works at millimetre
scale, not centimetre), the temperature sensor away from the
regulator's warm side, and the analog path --- SENSE to R1 to FILT
--- short and uncrossed.

\begin{figure}[htbp]
\centering
\begin{tikzpicture}[font=\scriptsize, >=latex]
% routed board, top view
\draw[thick, fill=graybg] (0,0.6) rectangle (9.0,4.4);
\node[font=\small\bfseries] at (4.5,4.15) {the routed 40$\times$30 board};
% header
\foreach \y in {1.2,1.9,2.6,3.3}{
  \draw[fill=white] (0.6,\y) circle (0.14);}
\node[left, font=\tiny] at (0.55,2.25) {J1};
% LDO + caps
\draw[fill=accentsoft] (2.6,1.1) rectangle (3.8,1.7);
\node[font=\tiny] at (3.2,1.4) {U1};
\draw[fill=white] (4.4,1.1) rectangle (4.9,1.5);
\draw[fill=white] (4.4,1.8) rectangle (4.9,2.2);
% sensor
\draw[fill=white] (7.4,2.6) circle (0.6);
\node[font=\tiny] at (7.4,2.6) {TMP36};
% RC
\draw[fill=white] (6.2,3.4) rectangle (6.8,3.8);
\node[font=\tiny] at (6.5,3.6) {R1};
\draw[fill=white] (7.0,3.4) rectangle (7.6,3.8);
\node[font=\tiny] at (7.3,3.6) {C1};
% tracks
\draw[thick, accent] (0.6,1.2) -- (2.6,1.2);
\draw[thick, accent] (3.8,1.4) -- (6.5,1.4) -- (6.5,3.4);
\draw[thick, accent] (7.4,3.2) -- (7.4,2.0) -- (8.4,2.0);
\draw[thick, good] (7.0,3.6) -- (6.8,3.6);
\node[align=left, text=slate, text width=10cm] at (4.6,-0.2)
{Placement rationale: header on the left edge (cable exits), the LDO
near the input with its decouplers tight, the sensor far from the
LDO's warm side, and the analog path short --- SENSE to R1 to FILT
in one clean run. The bottom face is one unbroken ground pour.};
\end{tikzpicture}

\caption{The routed board, top face: placement and the five tracks
that carry every net. Regenerate: \ttfamily
figures/ch18/fig\_routed.tex.}
\label{fig:routed18}
\end{figure}

\begin{example}[DRC, run on the real geometry]
\texttt{ch18\_drc\_check.py} applies the board's design rules ---
minimum clearance, minimum track width, hole keep-out --- to the
actual pad coordinates the Gerber generator emits. The real report:

\begin{lstlisting}[caption={\ttfamily ch18\_drc\_check.py output over
the real layout.}]
mount(2,2) pad-clearance 2.03mm
mount(38,28) pad-clearance 17.09mm
0 violations
\end{lstlisting}

Boring output is the point: a clean board's report should be a list
of facts that pass. To feel what it catches, move a mounting hole
0.5\,mm toward the header pads and re-run --- the clearance line
turns red the way KiCad's own DRC would.
\end{example}

\section{Fab files: the Gerber contract}\label{sec:gerber18}

\index{Gerber}
The deliverable a fab manufactures is a bundle: one file per copper
layer, one per mask and silkscreen face, one edge-cuts outline, and
an Excellon drill sheet. The chapter's \texttt{ch18\_gerber\_gen.py}
writes a real minimal set --- F\_Cu pads and the six tracks, B\_Cu
pour, Edge\_Cuts, and the drill file --- in valid Gerber X2 syntax:

\begin{lstlisting}[caption={\ttfamily ch18\_sensor-F\_Cu.gbr head ---
the header a Gerber viewer reads first.}]
%TF.FileFunction,Copper,L1,Top*%
%FSLAX34Y34*%
%MOMM*%
%LPD*%
%ADD10C,0.250*%
%ADD11R,1.800X1.800*%
%ADD12C,1.600*%
%ADD13O,1.300X2.300*%
\end{lstlisting}

\begin{figure}[htbp]
\centering
\begin{tikzpicture}[font=\scriptsize, >=latex]
% gerber -> fab -> assembly flow
\draw[thick, fill=accentsoft] (0.2,2.6) rectangle (2.6,4.2);
\node[align=center] at (1.4,3.7) {Gerber set};
\node[align=center, font=\tiny] at (1.4,3.1) {F\_Cu, B\_Cu, Edge\_Cuts,\\mask, paste, .drl};
\draw[->, thick] (2.6,3.4) -- (3.4,3.4);
\draw[thick, fill=graybg] (3.4,2.6) rectangle (5.8,4.2);
\node[align=center] at (4.6,3.7) {fab};
\node[align=center, font=\tiny] at (4.6,3.1) {etch, plate, mask,\\silk, test};
\draw[->, thick] (5.8,3.4) -- (6.6,3.4);
\draw[thick, fill=goodbg] (6.6,2.6) rectangle (9.0,4.2);
\node[align=center] at (7.8,3.7) {bare board};
\node[align=center, font=\tiny] at (7.8,3.1) {in the mail,\\$\sim$US\$5--15 + ship};
% checks box
\draw[thick, fill=amberbg] (3.0,0.6) rectangle (6.4,2.0);
\node[align=left, text width=3.2cm] at (4.7,1.3)
{\textbf{arrival checks}\\continuity on rails\\no shorts on pour\\pad sizes vs parts};
\draw[->, dashed, slate] (7.8,2.6) -- (6.4,1.4);
\node[align=left, text=slate, text width=10cm] at (4.7,-0.5)
{The Gerber is the contract the fab signs: every file is a layer or
a drill sheet. Our script-generated set is minimal but real ---
the fab's own DRC will run it again before the copper etches.};
\end{tikzpicture}

\caption{Gerber $\to$ fab $\to$ bare board $\to$ arrival checks.
Regenerate: \ttfamily figures/ch18/fig\_gerber.tex.}
\label{fig:gerber18}
\end{figure}

Ordering is a design decision too: $\sim$US\$5--15 plus shipping
gets five copies from the budget fabs, lead time a week or two ---
\emph{unverified, price-checked 2026-10-11}. The alternative is
keeping the board as a design exercise: every skill in this chapter
was exercised except the waiting.

When the board arrives: continuity-test the rails, check the pour
for shorts, hold the 0805s over their pads to confirm the footprints
matched the parts you bought.

\textbf{Safety recap.} A bought power supply, never a built one ---
that rule from ch1 and ch7 stands here too. If you solder this
board: iron in its stand between joints, lead-free solder, fan or
fresh air, wash hands. A LiPo pack powering the board gets the same
rules as ch17: protected cells, no shorts, no punctures.

\section{Common mistakes}\label{sec:cm18}

\begin{tripwire}[ERC warnings ignored]
ERC is free and runs in seconds; every ignored warning becomes a
\$15 lesson at the fab. Unconnected pins, power-output-pin warnings,
label typos --- clear them all, or write down why each one is fine.
\end{tripwire}
\begin{tripwire}[Schematic-footprint pinout mismatch]
The symbol's pin numbers must match the footprint's pad numbers ---
``pin 3'' on a TO-92 is not the same pin on an SOT-23. The netlist
carries names, not positions; assign the wrong footprint and the
board will politely build the wrong circuit.
\end{tripwire}
\begin{tripwire}[Missing decoupling at ICs]
The AMS1117 needs its capacitors \emph{at its pins} --- across the
board they are decoration. Every IC gets a local 100\,nF within a
few millimetres; the bulk cap guards the rail at entry.
\end{tripwire}
\begin{tripwire}[Trace width versus current]
A 0.15\,mm track is fine for a sensor signal and a fire hazard at
2\,A. Check current first, then width: our tracks are 0.25\,mm
because the whole board draws $\sim$10\,mA --- if it drew an amp,
the answer would be copper area, not a slightly wider line.
\end{tripwire}
\begin{tripwire}[Forgetting the ground return]
Signals routed neatly while ground wanders: the pour on the bottom
face exists so every return is short. On two-layer boards, ``route
signals, pour ground'' beats ``route everything.''
\end{tripwire}
\begin{tripwire}[Silkscreen as documentation]
Reference designators under parts, polarity marks hidden by the
component, no board name or date: the fab prints what you give it.
A board you can't identify next month is a board you'll reorder.
\end{tripwire}

\begin{exercise}[easy]
Why do we run ERC before assigning footprints?
\end{exercise}
\begin{exercise}[easy]
Name the five nets on this board and which layer each lives on.
\end{exercise}
\begin{exercise}[medium]
The fab charges per board size. Which single dimension change saves
the most cost on this layout without losing function?
\end{exercise}
\begin{exercise}[medium]
Sketch the keepout rule a mounting hole needs: what clearance to
copper, and what goes wrong if a track sneaks through?
\end{exercise}
\begin{exercise}[hard]
Argue for a copper-pour on the top face too --- when is a two-sided
pour the right answer, and what does it cost?
\end{exercise}

\section{Summary}\label{sec:sum18}

\begin{center}
\begin{tabularx}{\linewidth}{@{}lX@{}}
\toprule
\textbf{Thing} & \textbf{The rule of thumb}\\
\midrule
workflow & schematic $\to$ ERC $\to$ footprints $\to$ PCB $\to$ DRC
$\to$ Gerbers\\
schematic & the meaning lives here; fix bugs here, they're cheapest\\
footprint & value is schematic, body is footprint --- choose the
body you can solder\\
two layers & signals on top, unbroken ground pour on the bottom\\
DRC & boring reports are correct reports; clear every line\\
Gerbers & the fab's contract --- one file per manufactured layer\\
ordering & US\$5--15 + shipping, optional --- the skills transfer
either way\\
\bottomrule
\end{tabularx}
\end{center}

\section*{Additional exercises}
\addcontentsline{toc}{section}{Additional exercises}

\begin{exercise}[easy]
List the five Gerber-family files a minimal two-layer board needs.
\end{exercise}
\begin{exercise}[easy]
Why 0805 instead of 0402 for this book?
\end{exercise}
\begin{exercise}[easy]
Which three checks run when the board arrives, in order?
\end{exercise}
\begin{exercise}[easy]
What does the AMS1117's datasheet say the capacitors are for?
\end{exercise}
\begin{exercise}[medium]
Design the keepout for the TMP36's area: why give the sensor its own
quiet corner, away from the LDO?
\end{exercise}
\begin{exercise}[medium]
A track on the bottom face has to cross a signal --- on two layers,
what's the idiomatic fix that keeps the pour mostly intact?
\end{exercise}
\begin{exercise}[medium]
Estimate this board's power dissipation at the LDO and say whether
it needs thermal relief or a pour.
\end{exercise}
\begin{exercise}[medium]
The header pin-out must match the Tang Nano's GPIO row. Find which
J1 pins carry supply vs signal, and what happens if they're swapped.
\end{exercise}
\begin{exercise}[hard]
Add the mounting-hole keepout to the DRC rules: write the check as
it would appear in \texttt{ch18\_drc\_check.py} and test it on the
(2,2) hole.
\end{exercise}
\begin{exercise}[hard]
The 1-oz-copper question: at what current does a 0.25\,mm top track
stop being safe? Give the rule of thumb and the margin this board
has.
\end{exercise}
\begin{exercise}[hard]
Review a peer-style broken layout: TMP36 under the LDO's heatsink
side, C1 across the board, FILT track crossing three pads. List all
violations and the fix order.
\end{exercise}
\begin{exercise}[hard]
Extend the netlist to a second board revision adding an LED power
indicator on +3V3: what nets change, what parts, and which two ERC
checks does the change re-trigger?
\end{exercise}
\begin{exercise}[hard]
The pour-vs-routed-ground debate: measure (on paper) the return path
for U2's ground pin both ways and say why the pour wins for analog.
\end{exercise}
\begin{exercise}[hard]
Write the silkscreen policy a two-person team should adopt: what
must every board carry to be debuggable in a year?
\end{exercise}
\begin{exercise}[hard]
Cost the board end-to-end: fab + shipping + five sensors + passives.
Where does the budget break if you order twenty boards?
\end{exercise}
\begin{exercise}[hard]
The sensor feeds an FPGA threshold input. Describe the two honest
upgrades to a real ADC and what each costs in complexity.
\end{exercise}
\begin{exercise}[hard]
A ground loop appears when the board is powered by USB and probed
by a mains-referenced scope. Explain the mechanism and the safe
fix.
\end{exercise}
\begin{exercise}[hard]
Redesign the board on one layer: which net forces a jumper wire,
and what does that say about the two-layer convenience premium?
\end{exercise}

\section*{Hints and answers}
\addcontentsline{toc}{section}{Hints and answers}
\begin{description}
\item[ERC timing (easy, inline)] warnings cost seconds to fix on the
schematic and dollars on the board --- catch them where they're
free.
\item[Five nets (easy)] +5V, GND, +3V3, SENSE, FILT --- supply and
ground on the pour/fat tracks, the three signals on top tracks.
\item[Size change (medium)] shrinking 40\,mm to 25\,mm saves the
most area --- the header already forces the width.
\item[Mounting keepout (medium)] ~0.5\,mm copper-free ring around
the 3.2\,mm hole --- a track under a screw shorts to the standoff.
\item[Top pour (hard)] good for RF/EMI or extra return paths ---
costs etch time, solder-wick risk at pads, and one more place to
accidentally short.
\item[Gerber set (easy)] F\_Cu, B\_Cu, F/B Mask, Edge\_Cuts, drill
--- plus paste/silk for a nicer board.
\item[Why 0805 (easy)] hand-solderable by beginners and reworkable
with one iron; 0402 needs magnification and still hands.
\item[Arrival checks (easy)] rail continuity, pour shorts, parts-
over-pads fit.
\item[Capacitor purpose (easy)] LDO stability --- the datasheet's
minimum 22\,$\mu$F-ish guidance keeps the regulator out of
oscillation; decoupling also localises noise.
\item[Sensor keepout (medium)] the LDO's warmth biases a temperature
sensor; give TMP36 its own corner and the measured air wins.
\item[Bottom-face crossing (medium)] a short ``jumper track'' on the
top face lets the pour stay unbroken --- vias at both ends.
\item[LDO heat (medium)] $(5-3.3)\times$~10\,mA $\approx$ 17\,mW
--- trivially cool; the sensor's keepout is precaution, not need.
\item[Header swap (medium)] supply pins mis-wired: +5V into a
signal pin or a floating rail --- pin-number the silk and check
before powering.
\item[Keepout check (hard)] distance from (2,2) hole to every pad's
edge $\geq$ clearance --- implement as a second HOLES loop with
pad-radius subtraction.
\item[Track current (hard)] ~0.5--1\,A per 0.25\,mm at 1\,oz
(external layer, small rise); this board draws ~10\,mA --- margin
$\approx$ 100$\times$.
\item[Broken layout review (hard)] sensor near LDO (bias), C1
decoupling at distance (noise), FILT crossing pads (shorts) ---
fix order: placement, then pour, then route.
\item[Second revision (hard)] new parts: LED + resistor; net change:
+3V3 branches to LED-anode, LED-cathode$\to$res$\to$GND; ERC
re-checks power-pin connections and unconnected pins.
\item[Pour vs routed ground (hard)] U2's ground leg: pour path is
$\sim$2\,mm drop to sheet vs a centimetre-long routed track ---
lower impedance, less noise pickup on the analog reference.
\item[Silkscreen policy (hard)] board name+rev+date, refdes on every
part, polarity on every polar part, pin-1 marks, and a website/note
for who owns it.
\item[Twenty-board cost (hard)] fab price scales weakly per-board;
the sensor ($\sim$US\$2--3 each) becomes the dominant cost ---
budget breaks on the BOM, not the PCB.
\item[Real ADC upgrades (hard)] (1) a dedicated ADC IC (I2C/SPI)
--- one part, a bus, real code; (2) the FPGA's LVDS pins as a
comparator+SAR logic --- free pins, but Verilog calibration work.
\item[Ground loop (hard)] USB-ground and scope-ground share earth
potential through different paths --- current flows through your
measurement; fix: battery/isolated scope or lift one reference,
never defeat the safety earth.
\item[One-layer redesign (hard)] the GND return forces the jumper:
one wire over the top --- the 10-cent jumper is why two-layer
exists; the premium is cheap.
\end{description}

\section*{Further reading}
\addcontentsline{toc}{section}{Further reading}

\begin{itemize}
\item KiCad Documentation Team, \emph{Getting Started in KiCad}
(version-stamped 10.0.7).\newline
\url{docs.kicad.org/10.0/en/getting_started_in_kicad/} ---
the official walkthrough our pipeline follows.
Beginner--intermediate. Verified 200, 2026-10-11.
\item SparkFun, \emph{Beginner's Guide to KiCad}.\newline
\url{learn.sparkfun.com/tutorials/beginners-guide-to-kicad} ---
footprints, ERC/DRC habits, first-board ordering. Beginner.
Verified 200, 2026-10-11.
\item KiCad Documentation index.\newline
\url{docs.kicad.org/} --- schematic/PCB editor references for every
tool this chapter names. Reference. Verified 200, 2026-10-11.
\item JLCPCB.\newline
\url{jlcpcb.com} --- the budget-fab endpoint the Gerber bundle
targets; check its live quote for the current price. Verified 200,
2026-10-11.
\item Wikipedia, \emph{KiCad}.\newline
\url{en.wikipedia.org/wiki/KiCad} --- the project's history and
scope (why it's the free tool that won). Reference. Verified 200,
2026-10-11.
\item All About Circuits, \emph{Experiments} (vol.~VI) --- the build
exercises this board descends from. [unverified: 403 this session;
prior verified access]
\end{itemize}

\begin{rulebox}[{Watch instead}]
\begin{enumerate}
\item \emph{KiCad for Beginners --- Step by Step Tutorial (2025)}
--- DIY Hideout, 23:37,
\url{https://www.youtube.com/watch?v=d9_-lQq8ShE} --- the whole
pipeline this chapter walks, on screen. Matches \S\ref{sec:workflow18}.
Opened and confirmed 2026-10-11.
\item \emph{KiCad 6 STM32 PCB Design Full Tutorial} --- Phil's Lab
\#65, 1:40:05, \url{https://www.youtube.com/watch?v=aVUqaB0IMh4}
--- the deep version: a real MCU board start-to-finish including the
AMS1117-3.3 we use. Matches \S\ref{sec:spec18}--\S\ref{sec:gerber18}.
Opened and confirmed 2026-10-11.
\end{enumerate}
\end{rulebox}

\chapternotes{The pipeline and tool names follow the KiCad 10.0.7
docs; footprint choices and DRC habits follow SparkFun's guide;
the sensor-board schematic, netlist, DRC check, and Gerber bundle
are this book's own generated artifacts (all listed outputs are real
runs). Fab pricing is honest-marked and date-stamped. No KiCad GUI
was run in this session --- the design files are generated by the
chapter's own scripts and the KiCad steps are described-not-run.}
